% Registered A.48 after checkpoint07 release; full parent proof/source audit and independent certificate PASS.
\subsection{Normalizers in every two-power dihedral cover}

\begin{statement}{A.48}{Proved}{All-prime normalizers for the four-cusp two-power family}
For every \(m=2^s\), \(s\ge1\), every coefficient
prime \(\ell\), and the four-cusp curve \(Y_m\) of
A.44, the nontrivial closed procyclic subgroups of
\(\pi_1(Y_{m,\bar K})^{(\ell)}\) with open
arithmetic normalizer are exactly the nontrivial closed
subgroups of conjugates of cusp inertia. Together with
A.44 this proves both assertions of Task A.5 for the
entire two-power dihedral family. In particular it
settles the previously displayed case
\(v^2=w^{16}-1\) at \(\ell=17\).

The proof combines the actual compact-type
component-support obstruction of A.47 with a bound
on the multiplicity of a Galois character in the
primitive cyclotomic factor. It does not require an
explicit stable graph for these higher-genus curves.
\end{statement}

\begin{proof}
\paragraph{The primitive factor in the actual double-cover tower.}
Write \(B_m:v^2=w^{2m}-1\) and
\(\tau(w,v)=(\zeta_{2m}w,v)\). All constructions
may be defined over a finite extension containing the
indicated root of unity. There is an actual finite
etale map of punctured curves
\[
 q_m:Y_m\longrightarrow Y_{m/2},\qquad
                         (w,v)\longmapsto(w^2,v),
                    \qquad m\ge4.                  \tag{A.48.1}
\]
The corresponding proper map has degree two and
ramifies only at the four removed points over zero
and infinity. Its deck involution is \(\tau^m\).
The double-cover pullback and trace on rational
homology satisfy the usual identities
\(q_{m,*}q_m^*=2\) and
\(q_m^*q_{m,*}=1+\tau^m\). Consequently
\[
 W_m:=\ker(q_{m,*}:V_\ell J(B_m)\to V_\ell J(B_{m/2}))
                 =\{v:\tau^m v=-v\}.              \tag{A.48.2}
\]
This is the primitive \(\mathbf Q(\zeta_{2m})\)
factor: after adjoining the eigenvalues of \(\tau\),
it has exactly one eigenline for each odd exponent
modulo \(2m\). Every such eigenvalue has exact
order \(2m\). The eigenvalue assertion follows also
from the differential basis
\(w^{j-1}dw/v\), \(1\le j\le m-1\), and its
complex conjugate, as in A.44.

We emphasize that (A.48.1) is used covariantly on
fundamental groups and on homology. There is no
claim that an unramified double cover gives an
inclusion of pro-\(\ell\) groups when \(\ell\)
is odd.

\paragraph{A Galois character occurs at most twice in \(W_m\).}
Assume \(m\ge4\). We prove this assertion over every
finite number field over which \(\tau\) is defined,
and after every further finite extension. Enlarge
that field, if needed, to a finite Galois field \(L\)
containing \(\mathbf Q(\zeta_{2m})\). Enlarge the
coefficient field to a finite extension of
\(\mathbf Q_\ell\) containing all \(2m\)-th roots
of unity. Since Galois commutes with \(\tau\),
it acts on each of the one-dimensional \(\tau\)
eigenspaces by a character. No simplicity assertion
for the CM factor is required.

Here is the exact bound on equality of these
characters. It is convenient first to use the dual
etale cohomology, with its functorial action by
\(\tau^*\); passing to homology inverts Galois
characters and preserves their multiplicities.
Set
\[
 \Phi_m=\{1,3,5,\ldots,m-1\}
                \subset(\mathbf Z/2m)^\times.       \tag{A.48.3}
\]
The holomorphic eigenspaces of the primitive factor
are exactly those indexed by \(\Phi_m\).
For an embedding of \(L\) into
\(\overline{\mathbf Q}_\ell\) sending the chosen
\(\zeta_{2m}\) to \(\zeta_{2m}^a\), the
Hodge--Tate weight-one eigenspaces are indexed by
\(a\Phi_m\), and their complement has weight zero.
This follows directly from the displayed differential
basis and the canonical, functorial Hodge--Tate
comparison
\cite[Theorem 2.2.3, printed p.13]{fghBrinonConrad}.
All units \(a\bmod 2m\) occur, because \(L\)
contains the cyclotomic field. The comparison applies
at every place over \(\ell\), with its labelled
field embeddings. The weights remain unchanged on
finite extensions by
\cite[Theorem 2.4.6, printed p.20]{fghBrinonConrad}.

Thus, if the two global characters indexed by
\(j,k\in(\mathbf Z/2m)^\times\) are equal on
\(G_L\), then
\[
       j\in a\Phi_m\quad\Longleftrightarrow\quad
       k\in a\Phi_m\qquad\text{for every }a.
                                                               \tag{A.48.4}
\]
Equivalently \(kj^{-1}\) stabilizes \(\Phi_m\)
under multiplication. This is a necessary condition
for equality, which suffices for the bound; no
converse or assertion about finite-order twists is
needed.

The stabilizer is exactly
\[
       \operatorname{Stab}(\Phi_m)=\{1,m-1\}.       \tag{A.48.5}
\]
To prove this without a CM-type classification, put
\(z=\zeta_{2m}\). The finite geometric series gives
\[
               \sum_{j\in\Phi_m}z^j
                         =\frac{2z}{1-z^2}.
\]
If multiplication by \(a\) preserves \(\Phi_m\),
the same sum with \(z\) replaced by \(z^a\)
is equal to it. Hence
\(z^a-z^{-a}=z-z^{-1}\). For \(t=z^a\),
this factors as
\[
                         (t-z)(t+z^{-1})=0.
\]
Therefore \(a=1\) or \(a=m-1\pmod{2m}\).
Both do preserve \(\Phi_m\): for odd
\(1\le j<m\), multiplication by \(m-1\)
gives \(m-j\bmod 2m\). This proves (A.48.5).
Together with (A.48.4), it bounds the multiplicity
of every global Galois character in \(W_m\) by two,
even after any finite extension. In particular, if
\(0\ne v\in W_m\) spans a Galois-stable rational
line, then
\[
          \dim_{\mathbf Q_\ell}
              \operatorname{span}\{\tau^k v:k\in\mathbf Z\}
                              \le2.                \tag{A.48.6}
\]
Indeed all nonzero eigencomponents of \(v\) must
carry its single Galois character. Extending the
coefficient field for this argument does not change
the dimension of the rational span.

\paragraph{A wild cyclic action forces a large genus.}
Let \(Z\) be a smooth proper positive-genus curve
over an algebraically closed field of characteristic
two, with two-rank zero, and let a cyclic group of
order \(N=2^r\), \(r\ge1\), act faithfully on it.
Then
\[
                         g(Z)\ge\frac{(N-1)(N-2)}6.
                                                               \tag{A.48.7}
\]
For completeness, the Deuring--Shafarevich formula
\cite[formula (6), printed p.5; Lemma 3.2 and
Remark 3.1, printed pp.7--8]{fghGiuliettiKorchmaros}
shows that this action has one short orbit, consisting
of one fixed point. Directly, if \(f'\) is the
quotient two-rank and the short orbit sizes are
\(d_j\), then
\[
                    -1=N(f'-1)+\sum_j(N-d_j).
\]
It forces \(f'=0\). Every short orbit contributes
at least \(N/2\), so there is exactly one, of
size one. There is thus exactly one ramification
point, with full inertia group.

Let \(u_1,\ldots,u_r\) be its upper jumps.
By \cite[Theorem 1.1(c), manuscript p.2]{fghObusPries},
\(u_1\ge1\) and \(u_i\ge2u_{i-1}\), so
\(u_i\ge2^{i-1}\). Hilbert's different formula
and the upper-to-lower conversion in
\cite[Section 2.4, manuscript p.4]{fghObusPries}
give, after telescoping,
\[
 \delta=(N-1)+\sum_{i=1}^r2^{i-1}u_i
       \ge (N-1)+\sum_{i=1}^r4^{i-1}
        =\frac{N^2+3N-4}{3}.                       \tag{A.48.8}
\]
Riemann--Hurwitz, with nonnegative quotient genus,
gives \(2g(Z)-2\ge-2N+\delta\), proving
(A.48.7). The same argument includes a supersingular
elliptic component; it does not require genus at
least two in the Deuring--Shafarevich step.

\paragraph{The induction and its component-orbit contradiction.}
The assertion for \(m=2\) is A.40, and that for
\(m=4\) is the stable \(C_4\) case of A.14.
At coefficient prime two it holds for every \(m\)
by A.44 and A.46. Fix an odd \(\ell\), assume
\(m\ge8\), and suppose the assertion known for
\(m/2\).

Suppose a noncuspidal procyclic subgroup of the
geometric pro-\(\ell\) group of \(Y_m\) has open
arithmetic normalizer. A.11 gives a nonzero proper
homology vector \(v\). Over a finite field
extension contained in the normalizer's arithmetic
image, its rational line is Galois-stable. This
uses only the action on homology: geometric inner
automorphisms act trivially there, so no splitting
of the normalizer extension is assumed.

The actual map (A.48.1) and the induction hypothesis
force \(q_{m,*}v=0\). Otherwise the image procyclic
subgroup is nontrivial and has the same open
arithmetic-normalizer property; in the target it is
cuspidal and hence has zero proper homology, a
contradiction. Thus \(v\in W_m\), so (A.48.6)
applies.

At two choose a split stable model over a common
finite extension. A.44 makes its graph a tree, and
every positive-genus component has two-rank zero,
since the good special Jacobian is the product of
the component Jacobians and has two-rank zero. The
compact-type obstruction of A.47, with coefficient
\(\ell\ne2\), says that \(v\) is supported
on a single positive-genus component \(Z\).
The automorphism \(\tau\) extends to the stable
model and permutes the corresponding orthogonal
homology blocks. This is compatible with the actual
good-Jacobian specialization.

Let \(d\) be the length of the orbit of \(Z\)
under \(\tau\). The first \(d\) translates of
\(v\) are nonzero vectors in different homology
blocks, hence linearly independent. Formula
(A.48.6) gives \(d\le2\), and so \(d=1\)
or \(d=2\).

The restriction of \(\tau^d\) to \(Z\) has
exact order
\[
                              N=2m/d.              \tag{A.48.9}
\]
It has order dividing this number because \(\tau\)
has order \(2m\). Conversely every nonzero
\(\tau\)-eigencomponent of \(v\in W_m\)
has primitive \(2m\)-th eigenvalue, so
\((\tau^d)^k v=v\) is possible only when
\(N\mid k\). An identity automorphism of \(Z\)
would fix its whole homology block and hence \(v\).
This proves the exact order assertion without
assuming that \(\tau\) acts faithfully on every
component of the model.

All \(d\) components in this orbit have the same
genus. By (A.48.7), their contribution to the total
genus \(m-1\) is at least
\[
 \begin{cases}
  \displaystyle\frac{(2m-1)(2m-2)}6,&d=1,\\[4pt]
  \displaystyle\frac{(m-1)(m-2)}3,&d=2.
 \end{cases}                                                   \tag{A.48.10}
\]
Both exceed \(m-1\) when \(m\ge8\): their
differences are respectively
\(2(m-1)(m-2)/3\) and \((m-1)(m-5)/3\).
This contradiction proves the induction step.
It excludes all noncuspidal candidates at every odd
coefficient prime, including primes at which the
cyclotomic factor does not split over
\(\mathbf Q_\ell\).

Every nonzero closed subgroup of a cusp inertia
group has open arithmetic normalizer after making
the cusp rational, as in A.47. This proves the
classification. For the concrete previously open
case \(m=8\), an orbit of one would require a
component of genus at least \(35\), and an orbit
of two would require two components of genus at
least \(7\), whereas the total genus is \(7\).

The exact certificate \texttt{code/a\_two\_power\_normalizers.py} checks
the differential characters and primitive kernel,
the interval stabilizer and labelled weight-profile
multiplicities, the Hilbert--Herbrand telescoping
identity, and both genus contradictions. The proof
itself is uniform in \(m\).
\end{proof}
